\section{成立条件と実用上の課題}
\label{sec:hhl-assumptions-and-limitations}

HHLアルゴリズムは，行列の次元$N$に対して対数的に増加する量子ビット数で解状態を表現できる．
さらに，適切な入力モデルのもとでは，解に関する特定の期待値を次元$N$へ直接比例しない時間で評価できる可能性がある
\cite{hhl2009}．
しかし，この利点は任意の線形回帰問題で自動的に得られるものではない．
量子線形回帰全体を評価するには，次の条件と課題を考慮する必要がある．

\subsection{行列のエルミート性と可逆性}
\label{subsec:hermiticity-and-invertibility}

HHLアルゴリズムの基本的な定式化では，$A$がエルミートかつ可逆であることを仮定する．
線形回帰の正規方程式で用いる$X^{\mathsf{T}}X$は実対称行列であるため，エルミート性を満たす．
ただし，$X$が最大列階数を持たない場合は可逆でない．
この場合，線形従属な特徴量の除去，特異値の切り捨て，またはリッジ回帰による正則化が必要となる．

\subsection{疎性とハミルトニアンシミュレーション}
\label{subsec:sparsity-and-hamiltonian-simulation}

HHLアルゴリズムでは，$e^{\mathrm{i}At}$を効率よく実行できなければならない．
原型のアルゴリズムでは，各行の非零要素数が少ない疎行列と，
非零要素の位置および値へ効率よくアクセスできる入力方法を仮定している\cite{hhl2009}．
しかし，$X$が疎行列であっても，積$X^{\mathsf{T}}X$は密行列となる場合がある．
したがって，訓練データの疎性だけからHHLアルゴリズムの効率性を結論づけることはできず，
実際に使用する$A$の構造を調べる必要がある．

\subsection{状態準備の計算量}
\label{subsec:state-preparation-cost}

HHLアルゴリズムを開始するには，$\bm{b}$を振幅符号化した$\ket{b}$を準備する必要がある．
また，予測段階では新しい入力$\bm{x}_{*}$から$\ket{x_{*}}$を準備する必要がある．
第\ref{sec:data-encoding}節で述べたように，高次元ベクトルを少数量子ビットへ格納できても，
任意の振幅状態を作る回路が常に効率的とは限らない．
状態準備にベクトルの次元と同程度の時間が必要であれば，HHL部分の利点が相殺される可能性がある．

\subsection{条件数と精度}
\label{subsec:condition-number-and-precision}

条件数$\kappa(A)$が大きい場合，小さな固有値の逆数を高い精度で表現する必要がある．
そのため，位相レジスタの量子ビット数，ハミルトニアンシミュレーションの精度，制御回転の精度，
および補助量子ビットの測定回数が増加する．
特に，式\eqref{eq:normal-equation-condition-number}が示すように，
正規方程式では$X$の条件数が二乗される点に注意が必要である．
正則化は条件数を改善できる一方，元の最小二乗問題とは異なる解を与えるため，
数値安定性と統計的な偏りの間にトレードオフが生じる．

\subsection{出力の読み出し}
\label{subsec:solution-readout}

HHLアルゴリズムの出力$\ket{w}$から回帰係数$w_1,w_2,\ldots,w_d$をすべて古典的に読み出すには，
一般に多数の測定が必要となる．
そのため，HHLアルゴリズムは回帰係数の全成分を必要とする用途よりも，
式\eqref{eq:prediction-as-overlap}のように，解状態の内積や期待値だけを求める用途に適している
\cite{hhl2009,wiebe2012}．
量子線形回帰の出力形式を設計する際は，最終的に必要な情報が量子状態から効率よく評価できるかを先に確認する必要がある．

\subsection{量子優位性を評価する範囲}
\label{subsec:end-to-end-evaluation}

量子線形回帰の計算量は，HHLアルゴリズム単体だけでなく，次の処理を含めて評価する必要がある．

\begin{enumerate}[(1)]
  \item 訓練データから$A$と$\bm{b}$を構成する処理
  \item $A$へアクセスし，$e^{\mathrm{i}At}$を実行する処理
  \item $\ket{b}$および$\ket{x_{*}}$を準備する処理
  \item HHLアルゴリズムを成功させるための反復
  \item 内積や期待値を所望の精度で推定するための測定
\end{enumerate}

これらを含む端から端までの計算量が古典手法より小さい場合に初めて，
量子線形回帰としての計算上の利点を主張できる．
したがって，行列の次元だけでなく，疎性，条件数，状態準備，要求精度，出力形式を明示して比較することが重要である．
